\section{Introduction}\label{sec:introduction}

Viterbo's conjecture~\cite{Viterbo2000} states that among convex bodies of a
given volume in $\mathbb R^{2n}$ the Euclidean ball has the largest symplectic
capacity. We use the Ekeland--Hofer--Zehnder (EHZ) capacity. For a convex body
$K$ with smooth boundary, $c_{\rm EHZ}(K)$ is the least action
$\int_\gamma\lambda$ among closed characteristics $\gamma$ on $\partial K$,
that is, closed curves tangent to the line field $\ker(\omega_0|_{T\partial K})$;
here $\omega_0=\sum_j dq_j\wedge dp_j$ and $d\lambda=\omega_0$. For nonsmooth
convex bodies, polytopes in particular, the capacity is the limit over smooth
approximations. In dimension four the conjecture reads
$\operatorname{sys}(K)\le 1$ for the \emph{systolic ratio}
\[
 \operatorname{sys}(K)=\frac{c_{\rm EHZ}(K)^2}{2\operatorname{vol}_4(K)},
\]
which equals one on balls and is invariant under translations, rescalings and
linear symplectic maps.

Haim--Kislev and Ostrover disproved the conjecture with a Lagrangian
product~\cite{HaimKislevOstrover2024}. Write
$\mathbb R^4=\mathbb R^2_q\oplus\mathbb R^2_p$. For planar convex bodies
$K_q,K_p$, the Lagrangian product $K_q\times_L K_p$ is the set of
$(q,p)$ with $q\in K_q$ and $p\in K_p$; each factor lies in a Lagrangian plane.
Let $P_5$ be the regular pentagon of circumradius one centered at the origin with
a vertex at $(1,0)$, and $R(\theta)$ the rotation by $\theta$. The body
\[
 K_{\rm HKO}=P_5\times_L R(-\pi/2)P_5
 \qquad\text{has}\qquad
 \operatorname{sys}(K_{\rm HKO})=\frac{3+\sqrt5}{5}\approx1.047.
\]

This thesis starts from the counterexample and asks how much progress a fixed
toolbox makes on the questions around it. The toolbox consists of high-performance
implementations of capacity algorithms, proofs by exact computation,
standard data-science methods, computing time on the LICCA cluster of the
University of Augsburg, and AI agents. The questions came out of the work. For
example, random perturbations of $K_{\rm HKO}$ did not increase its systolic
ratio. This suggested that $K_{\rm HKO}$ is a local maximum, and for
perturbations that keep ten facets this is now the main theorem of the thesis.
The results therefore answer different questions and are presented side by
side rather than as one argument.

One of the earliest questions was whether standard data science finds new,
interesting bodies with systolic ratio above one. It did not. Bodies above one
arise near $K_{\rm HKO}$, since the systolic ratio is continuous, but random
sampling, adaptive sampling and local search started from random bodies all
stayed below one. The data did show a pattern: bodies whose two-dimensional
faces have small total symplectic area, relative to volume, tend to have large
systolic ratio. Looking for patterns in data has helped pure mathematics
before; in knot theory, relations noticed in tables of computed invariants
were later proved. Here the approach goes one step further: AI agents ran the
data science and also did large parts of the theory work. For instance, from
the suggestion to control a neighborhood of $K_{\rm HKO}$ by several upper
bounds on the systolic ratio, a GPT-5-series model found a configuration of
twenty-six such bounds that suffices, and wrote the exact SageMath
verification. AI agents are therefore part of the method, and the thesis
reports what they contributed.

The main results are the following. Every ten-facet polytope sufficiently close
to $K_{\rm HKO}$ has systolic ratio at most $(3+\sqrt5)/5$, with equality only
for images of $K_{\rm HKO}$ under translations, rescalings and linear
symplectic maps. The ten-facet space has another local maximum: a rational
product of a hexagon and a quadrilateral with systolic ratio one, which is
therefore not a new counterexample. No product of two independently linearly
deformed regular pentagons exceeds the systolic ratio of $K_{\rm HKO}$.
Symplectic ridge area has a strong inverse rank association with the systolic
ratio in sampled populations, but exact formulas for rotated pentagon and
triangle products order the two quantities in opposite ways. Behind these
results stand Haim--Kislev's finite capacity formula~\cite{HK2017}, its
first-order variation, and a second capacity algorithm built on the local
passage geometry of Chaidez and Hutchings~\cite{CH2021}.

The rest of this introduction first places $K_{\rm HKO}$ among capacities,
dynamics and convex geometry, then describes the results in more detail and
ends with a guide to the chapters.
